% =====================================================================
% GR-1i: The Classical Tests of Gravitation in Conscious Point Physics
% Companion i to GR-1 (Local Gravitation from SSV Shell Broadcast).
% Discharges OPEN-GR-TESTS-1: full geodesic derivations of the four
% classical tests on the GR-1c isotropic Schwarzschild metric.
% Version history: documentation_suite/changelog-GR-1i.md
% =====================================================================
\documentclass[12pt]{article}

\usepackage{amsmath,amssymb,amsthm}
\usepackage{geometry}
\usepackage{hyperref}
\usepackage{booktabs}
\usepackage{enumitem}
\usepackage{parskip}
\geometry{letterpaper, margin=1in}

% ---- corpus macros (matching GR-1 / GR-1c conventions) ----
\newcommand{\lP}{l_P}
\newcommand{\tP}{t_P}
\newcommand{\mP}{m_P}
\newcommand{\EP}{E_P}
\newcommand{\SSV}{\text{SSV}}
\newcommand{\PSR}{\text{PSR}}
\newcommand{\LSP}{\text{LSP}}
\newcommand{\rS}{r_S}
\newcommand{\riso}{r_{\mathrm{iso}}}
\newcommand{\vrho}{\varrho}

\newtheorem{theorem}{Theorem}[section]
\newtheorem{proposition}[theorem]{Proposition}
\newtheorem{remark}[theorem]{Remark}

\title{GR-1i: The Classical Tests of Gravitation\\
\large Perihelion precession, light deflection, Shapiro delay, and
gravitational redshift as geodesic consequences of the CPP
shell-broadcast metric\\
\normalsize Companion~i to GR-1 --- Conscious Point Physics gravitation
series --- Version 0.1}
\author{Thomas Lee Abshier, ND\\ Hyperphysics Institute}
\date{19 August 2026 --- V0.1 (CONV-029 revisions, 20 August 2026)}

\begin{document}
\maketitle

\begin{abstract}
GR-1c established --- conditional on the PSR constitutive form, whose
SR-1 grounding stands at W2 viability strength (GR-1 \S7) --- that the
CPP shell-broadcast mechanism yields \emph{exactly} the Schwarzschild
metric in isotropic coordinates.  This companion works out the standard
consequences: the full geodesic derivations of the four classical tests
of gravitation on that metric.  We derive the timelike orbit equation
and the perihelion advance ($42.99''$/century for Mercury), the null
geodesic and the light deflection at the solar limb ($1.75''$), the
round-trip radar excess (Shapiro) delay at superior conjunction
($\sim\!233\,\mu$s, Earth--Venus, grazing), and the static
gravitational redshift ($\Delta\nu/\nu = gh/c^2 = 2.46\times10^{-15}$
over 22.5\,m; GPS net $+38.5\,\mu$s/day including the special-relativistic
term).  The claim discipline is stated once and holds throughout:
because the metric is exactly Schwarzschild, every value is
\emph{identical to the GR prediction by construction}.  The classical
tests therefore discriminate CPP from Newtonian gravity --- decisively,
at the factor of two in deflection and at the entirety of the
perihelion and Shapiro signals --- but they do not, and cannot,
discriminate CPP from general relativity.  This paper adds no new
physics and no adjustable input; it is entry-criterion compliance,
written so that a reader asking ``does this theory pass the classical
tests?'' finds the derivations in one independently citable place.
Every number is machine-verified in the series verify script
(\texttt{series\_gravitation/code/3228\_classical\_tests\_verify.py},
8/8~PASS), which computes the two dynamical tests both in closed form
and by independent numeric Binet-equation geodesic integration.
\end{abstract}

\noindent\textbf{Keywords:} Conscious Point Physics; classical tests of
general relativity; perihelion precession; light deflection; Shapiro
delay; gravitational redshift; Schwarzschild metric; isotropic
coordinates; geodesic integration; Space Stress Vector.

% =====================================================================
\section*{Plain Language Summary}
% =====================================================================
Any proposed theory of gravity faces four famous entry exams.  Mercury's
orbit slowly swivels --- its closest-approach point creeps forward by 43
seconds of arc per century beyond what the other planets' tugs explain.
Starlight grazing the Sun bends by 1.75 seconds of arc --- twice the
value a naive Newtonian light-corpuscle calculation gives.  Radar
signals skimming the Sun on their way to Venus and back arrive about a
quarter of a millisecond late.  And clocks run slower the deeper they
sit in a gravitational well --- measurably so over a 22-metre tower, and
operationally so for every GPS satellite, whose clocks must be corrected
by 38 microseconds per day.

The parent paper GR-1 showed that Conscious Point Physics, via one
substrate response formula, reproduces the exact geometry general
relativity assigns to a static mass.  This companion takes the exam:
starting from that geometry and nothing else, it derives all four
effects step by step, and gets all four numbers.  Because the geometry
is \emph{exactly} the one general relativity uses, the answers agree
with general relativity exactly --- the honest statement, made
throughout, is that these tests separate CPP from Newton, not from
Einstein.  Passing them is a requirement, not a triumph; this paper
exists so that the requirement is met in writing.

\raggedright
\tableofcontents
\newpage

% =====================================================================
% SECTION 1: Introduction
% =====================================================================
\section{Introduction}
\label{sec:intro}

GR-1, the series parent, carries the predicted-versus-observed summary
of the four classical tests as a results-only table (GR-1
Table~1) and defers the derivations, by founder ruling, to one
dedicated companion: this paper.  The gap it closes is real ---
across the eight prior companions GR-1a--GR-1h, perihelion precession
and Shapiro delay had zero mentions and the word ``geodesic'' appeared
in one paper of eight --- and the closure is deliberately bounded:
standard geodesic integration on an already-published metric, with
every target number frozen in GR-1 Table~1 and machine-checked in the
series verify script before this paper was written.  The paper derives
what the script checks.

\subsection{Open Problems Addressed}
This paper discharges \texttt{OPEN-GR-TESTS-1} (the classical-tests
companion; registered Patch 3229, \texttt{frontier\_sectors/GR.md}).
It does not touch \texttt{OPEN-GR-FE-1} (derivation of the general CPP
field equations), which remains the deep open item of the arc; every
result here is a consequence of the recovered \emph{solution}, not of
any claimed field equation.

\subsection{Claim discipline (frozen)}
\label{sec:claims}
Three statements govern everything below.
\begin{enumerate}[leftmargin=2em]
\item \textbf{Conditionality.}  All results are conditional on the PSR
constitutive form $\PSR_{\rm eff} = \lP/(1 + k\,\Delta|\SSV|)$, whose
SR-1 grounding stands at W2 viability strength.  The inheritance is
stated in full in GR-1 \S7 and is not upgraded here.
\item \textbf{GR-identity by construction.}  Given that form, the
recovered metric is exactly Schwarzschild (GR-1c Theorem~1; identity
with the standard form verified at machine precision,
$3.3\times10^{-16}$).  Every classical-test value below is therefore
the GR value \emph{by construction}.  Nothing in this paper is a novel
prediction relative to GR.
\item \textbf{What the tests discriminate.}  The tests discriminate
CPP from Newtonian gravity (which fails deflection by a factor of two
and predicts no perihelion advance and no Shapiro delay at these
orders); they do not discriminate CPP from GR.  To pre-empt a
misreading (CONV-029, DeepSeek-seat note): CPP \emph{reproduces} GR's
classical tests; it does not \emph{share} GR's theoretical basis ---
the mechanism beneath (rigid GP lattice, compressed-DP aggregates, PSR
response) is entirely different, and the agreement is a derived
consequence, not a reformulation.  Frame-dragging
(Lense--Thirring) is treated in GR-1f and GR-1c and is not repeated
here.
\end{enumerate}

% =====================================================================
% SECTION 2: The metric
% =====================================================================
\section{The metric and its standard form}
\label{sec:metric}

The object under test is the GR-1c static solution.  With
$\vrho = GM/(2c^2\riso)$, the shell-broadcast source relation
$k\,\Delta|\SSV| = GM/rc^2$ and the PSR response yield
\begin{equation}
ds^2 \;=\; -\left(\frac{1-\vrho}{1+\vrho}\right)^{\!2} c^2\,dt^2
\;+\; (1+\vrho)^4 \left( d\riso^2 + \riso^2\, d\Omega^2 \right),
\label{eq:iso}
\end{equation}
exactly the Schwarzschild metric in isotropic coordinates.  Under the
exact coordinate transformation
\begin{equation}
r \;=\; \riso\,(1+\vrho)^2,
\label{eq:transform}
\end{equation}
Eq.~\eqref{eq:iso} takes the standard Schwarzschild form
\begin{equation}
ds^2 \;=\; -\Bigl(1 - \frac{\rS}{r}\Bigr) c^2\,dt^2
+ \Bigl(1 - \frac{\rS}{r}\Bigr)^{-1} dr^2 + r^2\,d\Omega^2,
\qquad \rS = \frac{2GM}{c^2}.
\label{eq:std}
\end{equation}
The identity of \eqref{eq:iso} with \eqref{eq:std} under
\eqref{eq:transform} is verified numerically at machine precision
(maximum relative $g_{tt}$ mismatch $3.3\times10^{-16}$ over five radii
spanning $10^9$\,m to 10\,AU; verify check~0).  Because the classical
tests are coordinate-invariant observables (proper-time ratios, asymptotic
deflection angles, round-trip proper-time delays measured by a distant
static clock), we are free to derive them in the standard coordinates
\eqref{eq:std}, where the geodesic equations take their most familiar
form; the results apply verbatim to \eqref{eq:iso}.  Section
\ref{sec:mechanism} returns to the isotropic form, where the CPP
mechanism is most transparent.

% =====================================================================
% SECTION 3: Geodesics
% =====================================================================
\section{Geodesics of the static metric}
\label{sec:geodesics}

Motion in \eqref{eq:std} conserves two quantities.  For an affinely
parametrized trajectory $x^\mu(\lambda)$ confined (without loss of
generality, by spherical symmetry) to the equatorial plane
$\theta = \pi/2$,
\begin{equation}
\tilde E \;=\; \Bigl(1-\frac{\rS}{r}\Bigr) c^2\,\frac{dt}{d\lambda},
\qquad
h \;=\; r^2\,\frac{d\varphi}{d\lambda}
\label{eq:conserved}
\end{equation}
(specific energy and specific angular momentum).  Substituting into
$g_{\mu\nu}\dot x^\mu \dot x^\nu = -\epsilon c^2$
($\epsilon = 1$ timelike with $\lambda = \tau$; $\epsilon = 0$ null)
gives the radial equation; with $u \equiv 1/r$ and
$'\equiv d/d\varphi$,
\begin{equation}
u'^2 \;=\; \frac{\tilde E^2}{c^2 h^2} - \bigl(\epsilon\,\tfrac{c^2}{h^2}
+ u^2\bigr)\bigl(1 - \rS u\bigr).
\end{equation}
Differentiating with respect to $\varphi$ yields the Binet equations
\begin{align}
\textbf{timelike:}\quad & u'' + u \;=\; \frac{GM}{h^2}
+ \frac{3GM}{c^2}\,u^2,
\label{eq:binet-m}\\[2pt]
\textbf{null:}\quad & u'' + u \;=\; \frac{3GM}{c^2}\,u^2.
\label{eq:binet-p}
\end{align}
The Newtonian limits are the same equations without the
$3GM u^2/c^2$ term: a closed Keplerian ellipse in the timelike case, a
straight line ($u = \sin\varphi/b$) in the null case.  Every classical
dynamical effect below is the work of that single relativistic term ---
which, in CPP language, is the first place the \emph{unexpanded}
nonlinearity of the PSR response becomes observable (GR-1 \S3).

% =====================================================================
% SECTION 4: Perihelion precession
% =====================================================================
\section{Test I: Perihelion precession of Mercury}
\label{sec:perihelion}

\subsection{Perturbative solution}
Write \eqref{eq:binet-m} as $u'' + u = A + B u^2$ with $A = GM/h^2$,
$B = 3GM/c^2$, and treat $B$ as small (for Mercury,
$B A \sim 10^{-7}$).  The zeroth-order solution is the Kepler ellipse
$u_0 = A(1 + e\cos\varphi)$.  Inserting $u_0$ into the perturbation and
retaining the resonant (secular) term $2A^2 B\,e\cos\varphi$ --- the
constant and $\cos 2\varphi$ pieces produce only bounded,
non-accumulating corrections --- gives
\begin{equation}
u \;\simeq\; A\bigl[1 + e\cos\bigl((1-\delta)\varphi\bigr)\bigr],
\qquad
\delta \;=\; A B \;=\; \frac{3\,(GM)^2}{c^2 h^2}
\;=\; \frac{3GM}{c^2\,a(1-e^2)},
\end{equation}
using $h^2 = GM\,a(1-e^2)$ for an orbit of semi-major axis $a$ and
eccentricity $e$.  Perihelion recurs when $(1-\delta)\varphi$ advances
by $2\pi$, i.e.\ after $\Delta\varphi = 2\pi/(1-\delta) \simeq 2\pi(1+\delta)$:
the orbit is an almost-closed ellipse whose apsidal line advances by
\begin{equation}
\Delta\varphi_{\rm prec} \;=\; \frac{6\pi GM}{a(1-e^2)\,c^2}
\quad\text{per orbit.}
\label{eq:prec}
\end{equation}

\subsection{Mercury}
With $a = 5.7909050\times10^{10}$\,m, $e = 0.205630$, orbital period
$87.9691$\,d, and $GM_\odot = 1.3275\times10^{20}$\,m$^3$\,s$^{-2}$,
Eq.~\eqref{eq:prec} gives
$\Delta\varphi_{\rm prec} = 5.019\times10^{-7}$\,rad per orbit.  Mercury
completes $36525/87.9691 = 415.2$ orbits per Julian century, so
\begin{equation}
\Delta\varphi_{\rm prec} \;=\; 42.99''\ \text{per century},
\end{equation}
against the observed anomalous advance $42.98 \pm 0.04''$/century
(the residual after planetary perturbations and solar oblateness are
removed).

\subsection{Independent numeric cross-check}
The verify script integrates \eqref{eq:binet-m} directly (RK4,
$d\varphi = 2\times10^{-5}$, 30 orbits, $\sim\!10^7$ steps), locates
perihelion passages as zero crossings of $u'$, and extracts the advance
by a telescoping estimator over 29 orbit intervals.  Numeric and
closed-form values agree to four decimal places:
$42.9917$ vs $42.9917$ arcsec/century.  Two numerical traps, found on
the script's first run and documented for any reimplementer, are
recorded in Section~\ref{sec:numerics}.

% =====================================================================
% SECTION 5: Light deflection
% =====================================================================
\section{Test II: Light deflection at the solar limb}
\label{sec:deflection}

\subsection{Perturbative solution}
For a photon with impact parameter $b$, the unperturbed solution of
\eqref{eq:binet-p} is the straight line $u_0 = \sin\varphi/b$
(incoming asymptote $\varphi = 0$, outgoing $\varphi = \pi$).  Writing
$u = u_0 + u_1$ and solving
$u_1'' + u_1 = (3GM/c^2)\sin^2\!\varphi/b^2$ gives
\begin{equation}
u_1 \;=\; \frac{3GM}{2c^2 b^2}\Bigl(1 + \tfrac{1}{3}\cos 2\varphi\Bigr)
\;=\; \frac{GM}{c^2 b^2}\bigl(1 + \cos^2\!\varphi\bigr)
\;\longrightarrow\;
\frac{2GM}{c^2 b^2}\ \text{at each asymptote.}
\end{equation}
Setting $u = 0$ at the outgoing asymptote $\varphi = \pi + \alpha_{1/2}$
and linearizing gives $\alpha_{1/2} = 2GM/(c^2 b)$ per asymptote, hence a
total deflection
\begin{equation}
\alpha \;=\; \frac{4GM}{c^2 b}.
\label{eq:defl}
\end{equation}
At the solar limb, $b = R_\odot = 6.957\times10^{8}$\,m:
\begin{equation}
\alpha \;=\; 1.75''
\qquad (1.7516''\ \text{in full precision}),
\end{equation}
confirmed by direct numeric integration of \eqref{eq:binet-p} to
$1.7517''$ (0.006\% agreement).  Observationally: the 1919 Eddington
expeditions at $\sim\!30\%$ precision, and modern VLBI at
$\gamma - 1 = (-0.8 \pm 1.2)\times10^{-4}$, i.e.\ the coefficient of
\eqref{eq:defl} confirmed to $10^{-4}$ relative precision.

\subsection{The factor of two}
\label{sec:factor2}
Half of \eqref{eq:defl} comes from time curvature ($g_{tt}$) and half
from spatial curvature ($g_{ij}$); a purely Newtonian (corpuscular or
time-curvature-only) treatment yields $2GM/(c^2 b) = 0.875''$.  In CPP
this split is structural, not accidental: the scalar broadcast
$|\SSV|_{\rm abs}$ sources $g_{tt}$ only, and GR-1b's extension of the
LSP broadcast to carry the net vector $\mathbf{\SSV}_{\rm net}$ ---
which sources $g_{ij}$ --- was forced by exactly this observable
(GR-1 \S4).  The classical deflection test is thus the single sharpest
discriminant in this paper: it separates CPP-with-two-component-LSP
from Newton \emph{and} from any scalar-broadcast-only variant of CPP
itself.

% =====================================================================
% SECTION 6: Shapiro delay
% =====================================================================
\section{Test III: Shapiro (radar echo) delay}
\label{sec:shapiro}

For a radar pulse traveling between coordinate radii $r_1$ and $r_2$
past the Sun with closest approach $b$, the coordinate time along the
null path of \eqref{eq:std} exceeds the flat-space value.  Integrating
$dt/dr$ along the null geodesic to first order in $\rS$,
\begin{equation}
\Delta t_{1\to2} \;=\; \frac{2GM}{c^3}
\ln\!\frac{4 r_1 r_2}{b^2}
\qquad \text{(leading logarithm, one way, } r_{1,2} \gg b\text{)},
\end{equation}
so the round-trip excess measured by a clock at Earth is
\begin{equation}
\Delta t \;=\; \frac{4GM}{c^3}\,\ln\!\frac{4\,r_E\,r_V}{b^2}.
\label{eq:shapiro}
\end{equation}
At Earth--Venus superior conjunction with a grazing ray
($r_E = 1$\,AU, $r_V = 1.0821\times10^{11}$\,m, $b = R_\odot$),
\begin{equation}
\Delta t \;=\; 232.6\,\mu\text{s} \;\approx\; 233\,\mu\text{s},
\end{equation}
the classic $\sim\!200\,\mu$s-class signal first measured by Shapiro's
group in 1968 at the few-percent level.  Sub-leading (non-logarithmic)
terms of order $GM/c^3$ contribute at the percent level and are
absorbed, in modern practice, into the parametrized coefficient
$\gamma$: the Cassini conjunction experiment constrains
$\gamma - 1 = (2.1 \pm 2.3)\times10^{-5}$.  Since the metric here is
exactly Schwarzschild, the PPN parameters are $\beta = \gamma = 1$
IDENTICALLY: compliance with the Cassini bound is a structural feature
of the formulation, not a numerical property of the leading-log term
--- which closes the loop on the sub-leading terms (CONV-029,
Gemini-seat note).  Accordingly, \eqref{eq:shapiro} at leading-log
precision is a COARSE consistency check, not a precision $\gamma$
test, and is presented as nothing more.

% =====================================================================
% SECTION 7: Gravitational redshift
% =====================================================================
\section{Test IV: Gravitational redshift and clock rates}
\label{sec:redshift}

\subsection{Static clock rates}
A clock at rest at radius $r$ in \eqref{eq:std} accumulates proper time
$d\tau = \sqrt{1 - \rS/r}\;dt$.  For emitter at $r_e$ and receiver at
$r_r$, the received-to-emitted frequency ratio is
\begin{equation}
\frac{\nu_r}{\nu_e}
\;=\; \sqrt{\frac{1 - \rS/r_e}{1 - \rS/r_r}}\,,
\qquad
\frac{\Delta\nu}{\nu}
\;\simeq\; \frac{GM}{c^2}\Bigl(\frac{1}{r_e} - \frac{1}{r_r}\Bigr)
\;\;\text{(weak field)}.
\label{eq:redshift}
\end{equation}
Over a small height difference $h$ near the Earth's surface this is the
uniform-field formula $\Delta\nu/\nu = gh/c^2$.  For the
Pound--Rebka--Snider 22.5\,m tower,
\begin{equation}
\frac{\Delta\nu}{\nu} \;=\; \frac{(9.80665)(22.5)}{c^2}
\;=\; 2.46\times10^{-15},
\end{equation}
measured to $\sim\!1\%$.  This test probes $g_{tt}$ alone --- in CPP
terms, the scalar $|\SSV|_{\rm abs}$ leg of the broadcast --- and is
the one classical test already implicit in the Newtonian-plus-clock-rate
layer of GR-1a.

\subsection{GPS: the operational test}
A GPS satellite ($r = R_E + 20{,}200$\,km) combines the gravitational
blueshift of \eqref{eq:redshift} with the special-relativistic time
dilation of its orbital speed $v = \sqrt{GM_E/r}$ (SR-1 physics,
inherited):
\begin{align}
\text{gravitational:}\quad
&+\frac{GM_E}{c^2}\Bigl(\frac{1}{R_E} - \frac{1}{r}\Bigr)
\times 86400\,\text{s} \;=\; +45.7\,\mu\text{s/day},\\
\text{special-relativistic:}\quad
&-\frac{v^2}{2c^2}\times 86400\,\text{s} \;=\; -7.2\,\mu\text{s/day},\\[2pt]
\text{net:}\quad &+38.5\,\mu\text{s/day}.
\end{align}
Uncorrected, this offset would accumulate a ranging error of
$\sim\!11$\,km/day; the correction is engineered into the satellite
clocks, making GPS a continuously operating confirmation of both terms.

% =====================================================================
% SECTION 8: Summary table
% =====================================================================
\section{Summary: predicted versus observed}
\label{sec:table}

Table~\ref{tab:tests} reproduces GR-1 Table~1; every number traces to
the derivations above and to the verify script.

\begin{table}[h]
\centering
\caption{Classical tests: CPP prediction (= GR value, since the metric
is exactly Schwarzschild) versus observation.  All values verified in
\texttt{series\_gravitation/code/3228\_classical\_tests\_verify.py}
(8/8 PASS).}
\label{tab:tests}
\begin{tabular}{lll}
\toprule
Test & CPP prediction & Observed \\
\midrule
Perihelion precession (Mercury) & $42.99''$/century &
$42.98 \pm 0.04''$/century \\
Light deflection (solar limb) & $1.75''$ &
$1.75''$ (Eddington 1919); \\
& & $\gamma - 1 = (-0.8\pm1.2)\times10^{-4}$ (VLBI) \\
Shapiro delay (superior conj.) & $\sim 233\,\mu$s (Earth--Venus,
grazing) & $\gamma - 1 = (2.1\pm2.3)\times10^{-5}$ (Cassini) \\
Gravitational redshift & $\Delta\nu/\nu = gh/c^2 = 2.46\times10^{-15}$
(22.5 m) & Pound--Rebka--Snider, $\sim1\%$ \\
 & GPS net $+38.5\,\mu$s/day & GPS operational \\
\bottomrule
\end{tabular}
\end{table}

The Newtonian column, were it added, would read: $0''$ (no relativistic
advance), $0.875''$ (half value), $0\,\mu$s (no coordinate-speed
retardation), and --- with clock rates grafted on via the equivalence of
potential and frequency shift --- the same redshift.  The discriminating
power of the table is entirely against that column.

% =====================================================================
% SECTION 9: Numerical verification and the two traps
% =====================================================================
\section{Numerical verification}
\label{sec:numerics}

The series verify script
(\texttt{series\_gravitation/code/3228\_classical\_tests\_verify.py})
performs, with no free parameters and no CPP-specific inputs beyond
``the metric is exactly Schwarzschild'': (0) the machine-precision
identity of \eqref{eq:iso} with \eqref{eq:std} under
\eqref{eq:transform}; (1) RK4 integration of \eqref{eq:binet-m} over 30
Mercury orbits, agreeing with \eqref{eq:prec} to $0.1\%$ (in practice,
four decimal places); (2) RK4 integration of \eqref{eq:binet-p},
agreeing with \eqref{eq:defl} to $0.006\%$; (3) evaluation of
\eqref{eq:shapiro}; (4) the Pound--Rebka and GPS numbers.  Result:
8/8 PASS.

Two numerical traps were found on the script's first run and are
documented here because either one silently corrupts the result at a
level large compared to the signal
(\texttt{series\_gravitation/reasoning/3228.md}):
\begin{enumerate}[leftmargin=2em]
\item \textbf{$\varphi$-accumulation drift.}  The integration angle
must be computed as $i\cdot d\varphi$ from an integer step counter.
Naive floating-point accumulation (\texttt{phi += dphi}) over
$\sim\!10^7$ steps injects $\sim\!10^{-6}$\,rad of rounding drift ---
percent-level against a perihelion signal of
$5\times10^{-7}$\,rad/orbit.
\item \textbf{Crossing overshoot.}  Zero crossings (perihelion
passages; the outgoing photon asymptote) must be located by linear
interpolation between bracketing steps.  Taking the first
post-crossing grid point overshoots by up to $d\varphi$ ---
10\%-level against a deflection of $8.5\times10^{-6}$\,rad.
\end{enumerate}

\subsection{Constants provenance and sensitivity}
\label{sec:constants}
Per the CONV-029 panel (constants-provenance question): this paper
deliberately quotes $GM_\odot = 1.3275\times10^{20}$\,m$^3$\,s$^{-2}$
--- the verify script's own $G \times M_\odot$ product --- so that the
paper's inputs ARE the script's inputs (provenance matching).  The IAU
nominal value is $GM_\odot = 1.3271244\times10^{20}$, a difference of
$0.028\%$.  Sensitivity of the table entries (all $\propto GM_\odot$
at leading order): perihelion $42.99 \to 42.98''$/cy (dead-centre of
the observed $42.98 \pm 0.04$), deflection $1.7516 \to 1.7511''$,
Shapiro $232.6 \to 232.5\,\mu$s; the redshift/GPS entries use Earth
constants and are unaffected.  Every shift is far inside the quoted
observational uncertainty; a reader preferring the IAU value may apply
the $0.028\%$ rescaling throughout.  Constants used: $G =
6.6743\times10^{-11}$, $c = 2.99792458\times10^{8}$, $R_\odot =
6.957\times10^{8}$\,m, $1\,\mathrm{AU} = 1.495979\times10^{11}$\,m
(as in the script).  The panel's independent SCRIPT-EXECUTED runs (two
seats; digits registered in \texttt{review/reviews-CONV-029.md})
serve as the archival verification artifact.  A caution the panel
endorsed: the numeric Binet integration agreeing with the closed forms
is an IMPLEMENTATION cross-check --- both routes instantiate the same
Schwarzschild geodesics --- not independent physical evidence; the
zero-new-predictions accounting of the Conclusion reflects exactly
this.

% =====================================================================
% SECTION 10: CPP physical mechanism
% =====================================================================
\section{CPP Physical Mechanism}
\label{sec:mechanism}

What is physically happening, in the substrate picture, during each
test?  The isotropic form \eqref{eq:iso} is the natural frame for the
answer, because in isotropic coordinates the spatial metric is
conformally flat and the two broadcast components appear separately:
$g_{tt} = -[(1-\vrho)/(1+\vrho)]^2$ carries the scalar
($|\SSV|_{\rm abs}$) content, and the conformal factor $(1+\vrho)^4$
carries the vector ($\mathbf{\SSV}_{\rm net}$) content.

\textbf{The medium.}  The source mass's polarization energy loads the
Dipole Sea with compressive stress, broadcast outward as shells of
$\LSP$ content.  The local response is the PSR shrinkage
$\PSR_{\rm eff} = \lP/(1 + k\,\Delta|\SSV|)$: Planck spheres nearer
the mass are smaller, so more processing steps fit per unit distance
--- the Sea is effectively \emph{denser} there.  Clocks (bound ZBW
cycles) tick slower in the denser region; light, which advances one
Planck sphere per Moment, has a lower coordinate speed there.  To first
order the coordinate light speed from \eqref{eq:iso} is
$c(r) \simeq c\,(1 - 2GM/c^2 r)$: the gravitating Sea behaves as a
graded-index medium with $n(r) \simeq 1 + 2GM/c^2 r$.

\textbf{The four tests, mechanistically.}
\begin{itemize}[leftmargin=2em]
\item \emph{Redshift} is the clock-rate gradient directly: the same
oscillator runs slower deep in the stressed Sea, so its emitted
frequency arrives low at a shallower receiver.  Pure scalar
($|\SSV|_{\rm abs}$) physics.
\item \emph{Light deflection} is refraction in the graded-index Sea:
wavefronts advance slower on the sun-side, turning the ray.  The
index $n(r) \simeq 1 + 2GM/c^2 r$ is frequency-independent, so the
bending is ACHROMATIC --- implicit in the metric's exactness, but
worth stating: chromatic dispersion of starlight at the limb would
falsify the mechanism as formulated (registered as an observation,
unminted; CONV-029, DeepSeek-seat note).  Exactly
half of $n - 1$ comes from the slowed clock rate (scalar) and half from
the compressed spatial lattice (vector); a scalar-only Sea would bend
light half as much (Section~\ref{sec:factor2}).
\item \emph{Shapiro delay} is the same slowed propagation integrated
along the path: the pulse spends extra coordinate time in the dense
region near the Sun, and a distant clock registers the accumulated
lateness.
\item \emph{Perihelion precession} is the orbit-scale signature of the
\emph{unexpanded} PSR nonlinearity: the effective potential acquires
the $\propto u^2$ correction of \eqref{eq:binet-m}, the apsides no
longer close, and the ellipse swivels.  Of the four tests it is the
only one sensitive beyond first order in $\vrho$ --- the first
observable to see that the PSR response is a full constitutive law
rather than a linearization.
\end{itemize}

\subsection{CP/GP Signature at This Scale}
\label{sec:signature}

\textbf{Load-bearing structure.}  The derivations rest on: the PSR
constitutive form (the arc's single nonlinear input, GR-1 Eq.~2;
conditional at W2 strength per the SR-1 inheritance), the
shell-broadcast source relation $k\,\Delta|\SSV| = GM/rc^2$ (exact,
GR-1a), and the two-component LSP broadcast (GR-1b), which supplies the
spatial half of the deflection and delay.  No axiom-level input enters
beyond what GR-1 already invokes; this paper adds only standard
geodesic mathematics.

\textbf{Visible versus smoothed discreteness.}  At solar-system scales
the lattice discreteness is entirely smoothed: the expansion parameter
$(\lP/r)^2$ is $\sim\!10^{-90}$ at Mercury's orbit, and the continuum
geodesic equations are exact to that accuracy.  Nothing in the four
tests can see a Planck sphere.  What remains visible is the
\emph{functional form} of the substrate response: the specific
$u^2$ correction in \eqref{eq:binet-m}, and the equal scalar/vector
split in the deflection, are fingerprints of the PSR formula and the
two-component broadcast respectively --- structure, not granularity.

\textbf{Macroscopic shadow.}  The classical tests are the macroscopic
shadow of the Dipole Sea's stress-response gradient: what conventional
physics describes as ``curved spacetime'' near a mass is, in CPP, the
smoothed manifestation of position-dependent Planck-sphere shrinkage
and directional Sea compression.  The factor of two in deflection is
the shadow of the broadcast's two-component structure.

% =====================================================================
% SECTION 11: Mapping to conventional physics
% =====================================================================
\section{CPP-to-Conventional-Physics Mapping}
\label{sec:mapping}

\begin{table}[h]
\centering
\caption{Structural correspondence for the classical tests.}
\label{tab:mapping}
\begin{tabular}{p{0.44\textwidth}p{0.48\textwidth}}
\toprule
CPP mechanistic element & Conventional description \\
\midrule
PSR shrinkage $\PSR_{\rm eff} = \lP/(1+k\,\Delta|\SSV|)$ &
Metric functions of the Schwarzschild solution \\
Scalar broadcast $|\SSV|_{\rm abs}$ (clock-rate gradient) &
$g_{tt}$; gravitational time dilation / redshift \\
Vector broadcast $\mathbf{\SSV}_{\rm net}$ (Sea compression) &
$g_{ij}$; spatial curvature; the $\gamma$-dependent half of
deflection and delay \\
Graded-index Sea, $n(r) \simeq 1 + 2GM/c^2 r$ &
Coordinate light speed in isotropic coordinates \\
Unexpanded PSR nonlinearity at orbital scale &
The $3GM u^2/c^2$ term of the relativistic Binet equation \\
Bound ZBW cycle ticking in the stressed Sea &
Ideal clock measuring proper time $d\tau = \sqrt{-g_{tt}}\,dt$ \\
\bottomrule
\end{tabular}
\end{table}

The correspondence is structural: CPP operates beneath the metric
description, and the tests confirm the smoothed (metric-level) output
of the mechanism, not the mechanism itself.  Any metric theory
producing exactly Schwarzschild produces this table's right column;
the left column is what CPP adds --- and what \texttt{OPEN-GR-FE-1}
must eventually ground in a derived field equation.

% =====================================================================
% SECTION 12: Conclusion
% =====================================================================
\section{Conclusion}
\label{sec:conclusion}

The four classical tests of gravitation follow from the CPP
shell-broadcast metric by standard geodesic integration, with no new
physics and no adjustable input, at values identical to general
relativity's by construction: $42.99''$/century, $1.75''$,
$\sim\!233\,\mu$s, and $gh/c^2$ with its GPS twin $+38.5\,\mu$s/day.
The entry criterion is met, in writing, in one citable place.  The
tests separate CPP from Newton --- completely --- and from
scalar-broadcast-only variants of CPP itself (via the deflection factor
of two); they do not separate CPP from GR, and the arc's genuinely
discriminating predictions remain where GR-1 locates them: the Planck
core and its echo/remnant consequences (GR-1d, GR-1e), which modify GR
where the classical tests cannot look.

\subsection{Swarm-Validation Contribution}
This paper mints \textbf{no new zero-parameter predictions}: by the
frozen claim discipline, its values are GR-identical consequences of an
already-published solution, and counting them as independent
correspondences would double-count GR-1c's exactness result.  Its
contribution to the swarm is qualitative --- entry-criterion
compliance: the programme's gravitational sector now passes, with
derivations on record, the four tests any gravitational theory must
pass before its distinctive predictions (the GR-1d/GR-1e echo and
remnant signatures, which \emph{are} counted in \texttt{predictions.md})
merit attention.

\subsection{Problem Status After This Paper}
\begin{itemize}[leftmargin=2em]
\item \texttt{OPEN-GR-TESTS-1}: OPEN $\rightarrow$ DISCHARGED at V0
(derivations complete; targets frozen in GR-1 Table~1 reproduced;
verify 8/8).  Final discharge at panel review / ship per the standard
cycle.
\item \texttt{OPEN-GR-FE-1}: unchanged (OPEN).  Nothing here bears on
the field-equation derivation.
\end{itemize}

\section*{Acknowledgements}
Derivations, drafting, and verification by the Anthropic Claude worker
under the CPP delegation protocol (PD-006), building on the GR-1c
solution and the Patch 3228 verify script.  The scope ruling (one
bounded tests companion; frozen targets; frozen claim discipline) is
the founder's (Session 148--149).

\begin{thebibliography}{9}

\bibitem{gr1} T.~L. Abshier, \emph{GR-1: Local Gravitation from SSV
Shell Broadcast}, CPP gravitation series, V0.2 (2026).  Parent paper;
Table~1 and the epistemic ledger govern this companion's claims.

\bibitem{gr1c} T.~L. Abshier, \emph{GR-1c: Strong-Field GR}
(formerly c08), CPP gravitation series (2026).  Theorem~1: the exact
isotropic Schwarzschild solution.

\bibitem{gr1b} T.~L. Abshier, \emph{GR-1b: Weak-Field GR}
(formerly c07), CPP gravitation series (2026).  The two-component LSP
broadcast and the deflection factor of two.

\bibitem{eddington} F.~W. Dyson, A.~S. Eddington, C.~Davidson,
``A Determination of the Deflection of Light by the Sun's
Gravitational Field,'' \emph{Phil. Trans. R. Soc. A} \textbf{220},
291--333 (1920).

\bibitem{shapiro64} I.~I. Shapiro, ``Fourth Test of General
Relativity,'' \emph{Phys. Rev. Lett.} \textbf{13}, 789--791 (1964);
I.~I. Shapiro \emph{et al.}, ``Fourth Test of General Relativity:
Preliminary Results,'' \emph{Phys. Rev. Lett.} \textbf{20}, 1265--1269
(1968).

\bibitem{poundrebka} R.~V. Pound and G.~A. Rebka,
``Apparent Weight of Photons,'' \emph{Phys. Rev. Lett.} \textbf{4},
337--341 (1960); R.~V. Pound and J.~L. Snider, \emph{Phys. Rev.}
\textbf{140}, B788 (1965).

\bibitem{cassini} B.~Bertotti, L.~Iess, P.~Tortora, ``A test of
general relativity using radio links with the Cassini spacecraft,''
\emph{Nature} \textbf{425}, 374--376 (2003).

\bibitem{vlbi} D.~E. Lebach \emph{et al.}, ``Measurement of the Solar
Gravitational Deflection of Radio Waves Using VLBI,''
\emph{Phys. Rev. Lett.} \textbf{75}, 1439 (1995); S.~B. Lambert and
C.~Le Poncin-Lafitte, \emph{Astron. Astrophys.} \textbf{529}, A70
(2011).

\bibitem{will} C.~M. Will, ``The Confrontation between General
Relativity and Experiment,'' \emph{Living Rev. Relativ.} \textbf{17},
4 (2014).

\end{thebibliography}

\end{document}
